% Section 8: inference (PHASE3)
\section{Inference}\label{sec:inference}

Source: PHASE3 (Appendix~\ref{app:phase3}). Checks: \path{verify_phase3.py},
17/17 PASS. Monte Carlo: \path{phase3_montecarlo.py}, 1000 replications per
design (Section~\ref{sec:evidence}).

\paragraph{Sampling (PHASE3 \S0.1).} \textbf{(P3-S1a)} Given the design and the
district effects, the unit rates are independent with $\E[\beta_i\mid\text{design}]=m(V_i)$.
\textbf{(P3-S2)} Inference is conditional on the design, so the sample-hull
vertices are fixed. The uncertainty is model-based, conditional on the design
(\S0.4, following Abadie, Athey, Imbens and Wooldridge~\citep{aaiw2023}): every unit is observed, nothing is assigned, and district shocks are fixed
features of this electorate (in the mean under M1/M2; a misspecification under
M0). No clustered variance is used. Independence is a substantive assumption: it
does not follow from observing every unit, and within-district dependence with
no mean shift violates it without being detected by the pooling test.
\begin{revnote}
(R-28, R-36; A2-5.) Round~1: the earlier P3-S1 allowed arbitrary within-district
dependence for M0 while P3.5(a) assumed independence; it was replaced by P3-S1a,
the \texttt{cluster} option was removed and the pooling test added. Round~2: the
two arguments (``assume independence given district effects'' versus ``all units
observed, so no dependence'') were separated; only the first is made.
\end{revnote}

\paragraph{Estimators (PHASE3 \S1).} OLS of each of $C-1$ columns on the
homogeneous quadratic monomials (plus district $\times X_r$ shifts for M1; one
regression per district for M2), with HC3 covariance. Each constraint is
$s_j(\beta)+a_j^\top k\ge0$ with $s_j$ linear in the coefficients $\beta$; the
target is $\tilde B(k)=c(\beta)+\gamma^\top k$ with sample weights. The plug-in
bound is $c(\hat\beta)+[\min,\max]\{\gamma^\top k:s(\hat\beta)+Ak\ge0\}$.

\begin{longtable}{@{}L{0.12\linewidth}L{0.13\linewidth}L{0.42\linewidth}L{0.25\linewidth}@{}}
\caption{Which result holds in which regime (PHASE3 \S0.5).}\label{tab:regimes}\\
\toprule Model & Regime & Formal result & Status \\ \midrule\endfirsthead
\toprule Model & Regime & Formal result & Status \\ \midrule\endhead
any rung & any, any $n$ & Theorem P3.8 (finite-sample CI and emptiness test) & Theorem; uninformative in width \\
M0 & any & P3.2, P3.4, P3.5(a), (b1), (b2) & Theorem (asymptotic) \\
M1, M2 & R1 ($D^\star$ fixed) & P3.2, P3.4, P3.5(a), (b1), (b2) & Theorem (asymptotic) \\
M1, M2 & R3 & the same with a high-dimensional Gaussian approximation & Conjecture \\
M1 & R2 & coverage of the CLR interval & Simulation finding \\
any & any & the CLR interval as implemented & Conjecture P3.5(b3) \\
\bottomrule
\end{longtable}

\begin{result}[Lemma P3.1, the endpoint is a minimum of linear functions --- PHASE3 \S2]\label{res:p31}
\status{Proved} (fixed constraint system). If the LP is feasible and bounded at
$\beta$, $\varphi_U(\beta)=c(\beta)+\min_{\lambda\in\Lambda}\lambda^\top s(\beta)$
with $\Lambda$ the finite set of dual vertices; $\varphi_U$ is concave and
piecewise linear in $\beta$.
\end{result}
\begin{plain}
Each end of the bound is the tightest of finitely many candidate limits, each
linear in the regression coefficients. JKST's case has four candidates and assumes
one wins clearly; in R$\times$C with districts there are many and near-ties are
the norm.
\end{plain}
\proofref[pE-res-lemma-P3.1]{Appendix~\ref{app:phase3}, \S2, steps 1--5}
\begin{revnote}(R-25.) Stated for a fixed constraint matrix and objective; joint variation in $(Q,\mu)$ is not covered.\end{revnote}

\begin{result}[Proposition P3.2, consistency and rate --- PHASE3 \S3]\label{res:p32}
\status{Proved.} Under P3-S1--S2, A1--A5, $\sqrt n(\hat\beta-\beta_0)=O_p(1)$ and
Slater (strict feasibility): (a) $\hat\varphi_U-\varphi_U(\beta_0;\hat A)=O_p(n^{-1/2})$;
(b) $\varphi_U(\beta_0;\hat A)\to\varphi_U(\beta_0;A)$ for nested or boundedly
convergent hulls; (c) since $\hat A_d\subseteq A_d$, estimating the hull only widens
the bound.
\end{result}
\begin{plain}
The plug-in bound converges at the usual rate to the bound for the sample hull,
which is weakly wider than the population bound. Coverage of the true parameter
never depends on the hull, because the true parameter satisfies the constraints at
every observed composition. Slater fails when an implied rate is exactly 0 at a
vertex (a structural zero); the rate result is then open, but coverage results do
not use Slater.
\end{plain}
\proofref[pE-res-proposition-P3.2]{Appendix~\ref{app:phase3}, \S3, steps 1--9; Remarks}
\begin{revnote}(R-26, R-32.) Non-nested hulls; Slater as an interiority condition; the error split into estimation ($O_p(n^{-1/2})$) and hull ($\ge0$, rate not shown) parts.\end{revnote}

\begin{result}[Theorem P3.4, asymptotic distribution for a fixed constraint system --- PHASE3 \S4]\label{res:p34}
\status{Cited structure} (directional delta method~\citep{fangsantos2019}).
If the constraint system is fixed along the sequence and
$\sqrt n(\hat\beta-\beta_0)\Rightarrow\mathcal Z\sim N(0,\Sigma)$, then
$\sqrt n(\hat\varphi_U-\varphi_U)\Rightarrow\min_{\lambda\in\Lambda^*}g_\lambda^\top\mathcal Z$.
(i) A single active candidate gives a normal limit. (ii) If two active gradients
differ in a direction where the estimator varies, the limit has strictly negative
mean: the plug-in bound is too narrow. (iii) Then the standard bootstrap is
inconsistent.
\end{result}
\begin{plain}
When several candidate limits are nearly tied, the plug-in end picks the one that
looks tightest by chance, so the estimated bound is biased inward. This is why the
plug-in is not a confidence interval.
\end{plain}
\proofref[pE-res-theorem-P3.4]{Appendix~\ref{app:phase3}, \S4, steps 1--6; Remark~P3.4.1 (use $C-1$ columns)}
\begin{revnote}
(R-27; A2-2.) Round~1: (ii) needs gradients that differ where $\Sigma$ has
variance (a null-direction difference gives a normal limit). Round~2: the theorem
was stated ``with the sample hull fixed'', which is not coherent as the hull grows;
it now requires a fixed constraint system (referee's $f_n=\min\{\theta,-\theta+n^{-1/2}\}$
gives $\min\{Z,1-Z\}$).
\end{revnote}

\begin{result}[Proposition P3.4$'$, triangular arrays --- PHASE3 \S4; register R17]\label{res:p34prime}
\status{Proved.} If the constraint system changes with $n$, $|\Lambda_n|\le\bar L$,
gradients converge and the localised slacks
$c_{\lambda,n}=\sqrt n(\theta_{\lambda,n}(\beta_0)-\varphi_{U,n}(\beta_0))\to c_\lambda\in[0,\infty]$,
then $\sqrt n(\hat\varphi_{U,n}-\varphi_{U,n}(\beta_0))\Rightarrow\min_{\lambda:c_\lambda<\infty}(g_\lambda^\top\mathcal Z+c_\lambda)$.
P3.4's limit obtains under separation ($c_\lambda\in\{0,\infty\}$).
\end{result}
\begin{plain}
A candidate that is inactive only by a gap of order $1/\sqrt n$ still shapes the
limit. The result describes the sample-hull endpoint; it says nothing about the
population-hull error.
\end{plain}
\begin{sketch}
By Lemma~P3.1 the scaled error equals exactly $\min_\lambda(g_{\lambda,n}^\top\sqrt
n(\hat\beta-\beta_0)+c_{\lambda,n})$; candidates with $c\to\infty$ drop out;
continuous mapping.
\end{sketch}
\proofref[pE-res-proposition-P3.4p]{Appendix~\ref{app:phase3}, \S4, Proposition~P3.4$'$ steps 1--4; check ``P3.4$'$ [F]''}
\begin{revnote}(A2-2.) New in round~2.\end{revnote}

\begin{result}[Theorem P3.5, confidence intervals --- PHASE3 \S5]\label{res:p35}
Fix levels $\alpha_U,\alpha_L,\alpha_E$; candidate upper endpoints
$\theta_\lambda(\beta)=g_\lambda^\top\beta$ with standard errors $\hat\sigma_\lambda$.
\begin{enumerate}[label=(\alph*),nosep]
\item \textbf{Conservative} (\status{proved}). Inflate every constraint and the target
constant by a simultaneous Gaussian critical value $\hat k$, solve the LP over the
inflated set $\hat\Theta^+$, and add $\pm z_{\alpha_E}S_1$ with
$S_1=\tfrac12(\sum_iw_{ir}^2)^{1/2}$. Then
$\liminf P(B_{rc}\in\mathrm{CI})\ge1-\alpha_U-\alpha_L-\alpha_E$; without $S_1$ it
covers $\tilde B$.
\item[(b1)] \textbf{CLR with the full candidate set} (\status{cited}: Chernozhukov, Lee
and Rosen~\citep{clr2013}, Thm~4, under Condition~P). Same coverage.
\item[(b2)] \textbf{Any candidate set with a simultaneous critical value}
(\status{proved}). For any data-dependent set of dual-feasible points,
$\hat U=\min_{\lambda\in\hat\Lambda}(\hat\theta_\lambda+\hat k\hat\sigma_\lambda)$
covers $\varphi_U$ with probability $\ge1-\alpha_U$ asymptotically if (i) $\hat k$ is
computed over a fixed finite superset, or (ii) $\hat k=\sqrt{\chi^2_{q,1-\alpha}}$
(projection) with $\Sigma$ positive definite and $q=\dim\beta$, or an exact reduced
full-rank parameterisation. An exact reduction means documented identities
(for example, adding-up removes column $C$); a rank taken from a numerical cutoff
on $\hat\Sigma$ is not one, so then $q=\dim\beta$ must be used, and dropping
candidates that load on numerically null directions is only a safeguard (it can
only widen the interval).
\item[(b3)] \textbf{The CLR procedure as implemented} (\status{conjecture /
algorithm}). Candidates are collected from the optimal dual vertices at
$\hat\beta$ and at simulated draws. A reduction to (b1) is proved when the
collected set contains CLR's selected set and a minimiser; recovery of rare
vertices, the downward effect on $k_p$, and growing $|\Lambda|$ are open.
\item[(c1)] \textbf{Variance-adaptive allowance} (\status{proved}; register R18).
Replace $S_1$ by the Bhatia--Davis~\citep{bhatiadavis2000} bound $\hat S_{BD}$ computed from a P3.5(a) or
projection interval for $\bar m=\sum_i\pi_im_i$ on the rung's own constraints;
if $\max_iw_i^2/\bar s^2\to0$, coverage is $\ge1-\alpha_{\rm est}-\alpha_E-\alpha_v$.
\item[(c2)] With a CLR interval for $\bar m$: coverage conditional on that interval
covering; inherits (b3)'s status.
\end{enumerate}
\end{result}
\begin{plain}
The confidence interval widens each candidate limit by its own uncertainty before
taking the tightest, so noisy candidates (for example, districts with few polls)
rarely set the limit. Three versions carry proofs: the conservative interval (all
constraints widened at once), the projection interval (one ellipsoid for the
whole coefficient vector; valid for any set of candidates, but wide when there
are many coefficients) and the fixed-superset version. The package's default,
CLR with simulated candidates, is narrower in grouped designs and covered in every
simulation design, but its validity is not proved: it is an \emph{algorithm}.
For the actual rate $B$, an extra allowance covers each unit's local luck; a
tighter, variance-adaptive allowance is valid if its own input interval comes
from a proved method.
\end{plain}
\begin{sketch}
(a) On one simultaneous event the true $k_0$ is feasible for the inflated
constraints, so $\tilde B$ is inside the inflated LP range; the residual term is
handled by Chebyshev or Lindeberg along subsequences. (b2) By weak duality every
dual-feasible candidate is an upper bound for $\varphi_U(\beta_0)$; a simultaneous
event for all candidates (a fixed superset, or Cauchy--Schwarz in the metric of
$\hat\Sigma$ for projection) then covers. (c1) Concavity of $m(1-m)$ bounds the
residual variance by $\bar s^2$; on the event that the $\bar m$ interval covers,
$\hat S_{BD}\ge\bar s$; a union bound finishes.
\end{sketch}
\proofref[pE-res-theorem-P3.5]{Appendix~\ref{app:phase3}, \S5: (a) steps 1--7; (b2) steps 1--4; (b3) proved reduction; (c1) steps 1--5}
\begin{revnote}
(R-29; P-B40, P-B41; I.3; A2-3.) Round~1 (Phases~1--3 referee): the simulated-candidate
CLR was claimed as a theorem; P3.5(b) was split into (b1) cited, (b2) proved and
(b3) conjecture, and the package labels CLR algorithmic with projection alongside.
Batch~B: (b2)(ii) took $q=\operatorname{rank}\Sigma$ with a pseudoinverse, which
does not control errors in null directions (counterexample $(Z/\sqrt n,-1/n)$);
it now needs $\Sigma$ positive definite or an exact reduction. The residual step
of (a) applied a CLT under the weight condition alone; the referee's
Chebyshev/Lindeberg subsequence argument replaced it. Round~2: a numerical rank
cutoff is a safeguard, not case (ii-b), so the package uses $q=\dim\beta$; the
projection interval for $B$ was not covered when its Bhatia--Davis bound came from
CLR, which led to P3.5(c).
\end{revnote}

\begin{result}[Proposition P3.6, tests --- PHASE3 \S5]\label{res:p36}
(a) \textbf{Emptiness test} (\status{proved}): reject the rung if $\hat\Theta^+$ is
empty; level $\le\alpha_U+\alpha_L$. (b) \textbf{M1 against M2}: Wald test of the
$(D^\star-1)(C-1)R(R-1)/2$ overidentifying restrictions; $\chi^2$ calibration
\status{proved for fixed dimension} (R1). (c) \textbf{Pooling, M0 against M1}
(\status{proved}): conditional on the design, with $R$ affinely independent
compositions per district, M0's mean model holds iff all district shifts vanish;
test $\zeta=0$ with $(D^\star-1)(C-1)R$ restrictions. It tests mean shifts only.
\end{result}
\begin{plain}
An empty confidence set is a real rejection at the stated level. The curvature
test checks M1's testable part. The pooling test says whether the pooled model is
wrong for this design; rejection does not mean its intervals fail (design S1), and
the test has no power against within-district dependence.
\end{plain}
\begin{sketch}
(a) On the simultaneous event of P3.5(a) the true $k_0$ lies in the inflated set,
so emptiness requires that event to fail. (c) If M0 held for the observed units,
then in each district $(G_c-\zeta_{c,d}\mathbf 1^\top)X_i=0$ for all its units, with
$G=H'-H$. Because the district has $R$ affinely independent compositions, which
lie in $\{\mathbf 1^\top x=1\}$ and are therefore linearly independent, this forces
$G_c=\zeta_{c,d}\mathbf 1^\top$ in every district. With $\zeta_{c,1}=0$, $G_c=0$ and all
shifts vanish. So M0 holds iff $\zeta=0$, a linear hypothesis on the M1
regression.
\end{sketch}
\proofref[pE-res-proposition-P3.6]{Appendix~\ref{app:phase3}, \S5, Proposition~P3.6 proofs of (a) and (c)}

\paragraph{Calibration by restricted wild bootstrap (package 0.3.0; PHASE3 \S5).}
\status{Simulation-supported, not proved for growing dimension.} In moderate
dimension the HC3 Wald statistic against $\chi^2$ over-rejects (pooling test, 38
restrictions: 0.078 at nominal 5\%). The package fits the null model, draws
$Y^*=\hat Y_0+\hat e_0v$ with Rademacher $v_i$ shared across the $C-1$ equations,
refits the larger model and recomputes the HC Wald statistic;
$p=(1+\#\{W^*\ge W\})/(B+1)$, $B=399$ by default. Rejection rates at 5\% in design
S0: pooling 0.043 (SE 0.005; $\chi^2$ 0.078, F 0.072); M1 vs M2 0.054 (SE 0.007;
$\chi^2$ 0.057). Scripts: \path{benchmarks/pooling_test_size.R},
\path{benchmarks/m1_test_size.R}.
\begin{revnote}(R-30, R-37; WS6; A2-5.) Degrees of freedom corrected; bootstrap added after the over-rejection was found; the pooling test is stated as a test of mean shifts only.\end{revnote}

\begin{result}[Theorem P3.8, finite-sample CI and emptiness test --- PHASE3 \S5.1]\label{res:p38}
\status{Proved.} Under P3-S1a, a correctly specified mean model and a full-rank
design, with Hoeffding~\citep{hoeffding1963} radii for every constraint and the target constant (union
bound over $J+2$ events) and $t_E=(\tfrac12\log(2/\alpha_E)\sum_iw_{ir}^2)^{1/2}$:
for every $n$ and in every regime, (a) $P(B\in\mathrm{CI}^{\rm fin}\mid\text{design})\ge1-\alpha_s-\alpha_E$;
(b) without $t_E$ it covers $\tilde B$ at $1-\alpha_s$; (c) the emptiness test has
exact level $\le\alpha_s$; (d) intersecting with DD keeps (a).
\end{result}
\begin{plain}
This interval needs no large-sample theory and no fixed number of districts, so it
is the only formal guarantee with many small districts. In every design tried it is
as wide as the logical range $[0,1]$, mainly because of the
constraint-by-constraint union bound. It is valid but not yet useful, and is not
in the package.
\end{plain}
\begin{sketch}
Every constraint term and the target constant are linear functionals of the OLS
fit, and each is a sum over units of independent terms whose range is known from
the design (the outcome shares lie in the simplex). OLS is unbiased under the
mean model, so Hoeffding's inequality gives a finite-sample radius for each. A
union bound over the $J+2$ one-sided events puts the true $k_0$ in the widened
set; the residual term gets its own Hoeffding radius. No Gaussian approximation
is used, so nothing depends on the regime.
\end{sketch}
\proofref[pE-res-theorem-P3.8]{Appendix~\ref{app:phase3}, \S5.1, steps 1--7}
\begin{revnote}(R-31; A2-6.) New after round~1; round~2 agreed it is a fallback guarantee and left informative R2 inference open.\end{revnote}

\begin{result}[Proposition P3.7, many districts with few units --- PHASE3 \S6]
\status{Proved.} (a) Dropping districts' constraints is conservative. (b) If the
retained districts are chosen from the design only, any procedure keeps its
conditional coverage; asymptotic results transfer only if their regime conditions
hold for the retained districts.
\end{result}
\begin{plain}
With many small districts, ignoring some districts' constraints can only widen
the bound, so it never costs validity of the identified set. If the retained
districts are chosen from the design alone (for example, those with at least 30
polls or enough spread in composition), not from the outcomes, any interval that
is valid for a fixed set of districts stays valid. Selection does not, by itself,
move the problem into a regime where the asymptotic theory applies.
\end{plain}
\proofref[pE-res-proposition-P3.7]{Appendix~\ref{app:phase3}, \S6, steps 1--2}
\begin{revnote}(R-31.) (b) was weakened from ``every result holds conditionally'' to a conditional statement.\end{revnote}

\begin{result}[Proposition S1, bounded nonlinearity --- PHASE3 \S7]\label{res:s1}
\status{Proved} (an outer bound). If $\E[\beta\mid V]=M_H(X)+\Delta(X)$ with $\Delta$ in
the undetectable class of PHASE2a Prop.~9 and $\|\Delta\|_\infty\le\eta$, then
$\tilde B_{rc}\in[\min_{\ThO(\eta)}\varphi-\eta,\max_{\ThO(\eta)}\varphi+\eta]$ with
$\ThO(\eta)=\{k:s+Ak\ge-\eta\mathbf 1\}$.
\end{result}
\begin{plain}
If rates may depart from straight lines, invisibly, by up to $\eta$, the bound
widens in a controlled way. The breakdown value $\eta^*$ of a conclusion is the
largest departure under which it survives; in the $3\times3$ example,
$\eta^*(\text{loyalty}>0.30)=0.060$. The same logic applies to the dials $\delta$
(district heterogeneity, change between periods).
\end{plain}
\proofref[pE-res-proposition-S1]{Appendix~\ref{app:phase3}, \S7, steps 1--3}
\begin{revnote}(R-34.) Labelled an outer bound under bounded deviations, not the sharp set of the relaxed model.\end{revnote}

\paragraph{Boundary cases and JKST's $\mathrm{CI}_x$ (PHASE3 \S\S4--5).} When the
true $K$ is on the boundary, coverage is close to nominal; in the interior it tends
to one (Imbens--Manski). Near point identification, one-sided intervals per
endpoint under-cover unless the critical value adapts~\citep{stoye2009}; the
two-sided split used here is uniformly valid at some cost in width. JKST's
$\mathrm{CI}_x$ is version (b) with a single candidate per side.

\paragraph{Practical recipe (PHASE3 \S9).} Fit on $C-1$ columns with HC3; report
the plug-in as an estimate of the identified set; report CLR (algorithmic) for
$\tilde B$ and $B$ with the projection interval alongside (for $B$, with $\bar m$
from the projection interval on the rung alone, or Popoviciu); intersect with DD;
report the emptiness, M1 and pooling tests; report sensitivity in $\eta$ and
$\delta$ with breakdown values; with few units per district, prefer selection,
M1($\delta$) or coarser groups, and say that CLR's coverage there rests on
simulation.
